\Needspace{15\baselineskip}
\fancyhead[L]{\small GAUSS AN VON ZACH.}
\begin{center}
GAUSS AN VON ZACH.
\end{center}

\begin{center}
Monatliche Correspondenz. Band XIV. S. 377–382. Oktober 1806.
\end{center}

\begin{flushright}
Braunschweig, 25. August 1806.
\end{flushright}

\par\indent --- --- --- Anbei erhalten Sie die Ephemeride für Juno. Sie ist zwar schon seit einiger Zeit fertig, doch ich hatte immer gehofft, die für \emph{Ceres} zugleich beifügen zu können. Von Ceres habe ich aber noch immer keine ganz zuverlässigen Beobachtungen aus der Zeit um die
\SourceLeaf{285}%
Opposition. --- Sollte ich nicht bald gute Beobachtungen von ihr erhalten, werde ich die Elemente vorerst mit den Beobachtungen von \textsc{Pasquich} verbessern müssen, die nach seinem eigenen Urteil nicht sehr zuverlässig sind. Die Beobachtungen von \textsc{Harding} und ein paar, die mir \textsc{Bode} mitgeteilt hat, liegen zeitlich schon etwas weiter von der Opposition entfernt. Sollte man sie nicht auch in Mailand beobachtet haben?

\par\indent Die Elemente, nach denen die Ephemeride der \emph{Juno} berechnet wurde, sind folgende:

\begin{center}
VI. \emph{Elemente der Juno.}
\end{center}

\Needspace{18\baselineskip}
\begin{center}
Epoche der Länge im Meridian von Seeberg
\end{center}
\[
\left\{\begin{array}{ll}
1805. & 42^\circ\,35'\,7''3\\
1806. & 125\;11\;20.1\\
1807. & 207\;47\;33.0
\end{array}\right.
\]
\begin{center}\begingroup
\fontsize{9}{11}\selectfont
\setlength{\tabcolsep}{3pt}\renewcommand{\arraystretch}{1.12}
\begin{tabular}{ll}
tägliche mittlere tropische Bewegung  & $814''7201$ \\
Länge des Aphels 1806 .\ .\ . & $233^\circ\,17'\,1''$ \\
Länge des aufsteigenden Knotens 1806  & $171^\circ\,4'\,57''4$ \\
Exzentrizität .\ .\ .\ .\ .\ .\ .\ .\ . & $0.254944$ \\
Logarithmus der Halbachse .\ .\ .\ . & $0.4260480$ \\
Neigung der Bahn .\ .\ .\ .\ .\ .\ . & $13^\circ\,3'\,28''4$ \\
\end{tabular}\endgroup\end{center}

\par\indent Prof. \textsc{Harding} schreibt mir, er habe beide Karten schon größtenteils fertiggestellt. Eine Unpässlichkeit verzögere nur noch den Vergleich mit dem Himmel selbst. Er hofft, sie Ihnen bald zuschicken zu können.

\Needspace{15\baselineskip}
\par\indent Die Lichtstärke der \emph{Juno} im nächsten Jahr habe ich in meiner Ephemeride nicht angegeben, weil ich sie diesmal nach einer neuen, besonders bequemen Methode berechnet habe, bei der der Abstand des Planeten von der Sonne überhaupt nicht in die Rechnung eingeht. Zur Übersicht gebe ich aber doch drei berechnete Werte für die Lichtstärke an, aus denen hervorgeht, wie außerordentlich schwierig die Beobachtung der Juno im nächsten Jahr sein wird:

\begin{center}\begingroup
\fontsize{9}{11}\selectfont
\setlength{\tabcolsep}{3pt}\renewcommand{\arraystretch}{1.12}
\begin{tabular}{ll}
1807 Januar 12 .\ . & $0.00661,$ \\
Mai 18 .\ . & $0.01600,$ \\
Sept. 21 .\ . & $0.00709$ \\
\end{tabular}\endgroup\end{center}

\par\indent In diesem Jahr betrug die größte Lichtstärke im März immerhin noch $0.03983$ (M.\,C., Band~XI, \href{https://www.digitale-sammlungen.de/en/view/bsb10538603?page=492}{S.~478} [\hyperlink{native-270}{S.~265} dieses Bandes]). 1807 wird sie also, wenn sie am größten ist, nicht einmal halb so groß sein. --- Die eben erwähnte neue Methode verkürzt die Rechnung so sehr, dass ich auf einer gewöhnlichen Oktavseite 15 vollständige Rechnungen unterbringe, wenn die Positionen der Sonne als bereits gegeben gelten.

\par\indent [Der geozentrische Lauf der Juno vom 12. Januar bis zum 21. September 1807 ist in Abständen von drei Tagen für 12 Uhr, Mitternacht in Seeberg, dargestellt. Rektaszension und Deklination sind auf die Minute, der Abstand von der Erde auf drei Dezimalstellen genau angegeben.]

